


\todo{neu mit SCC-NEA, ne geht nicht, in diskussion, ab links, aa rechts und sonst b, oben in zwei zustände aufsplitten,
nea idea nicht überträgbar von kommutativen}



\todo{erwähnen, dass ich diese konstruktion noch nie gesehen, und deswegen keine referenz...
in diskussion?}

\subsubsection{The State Merging NFA}

\todo{wird nur hier gebraucht, also nur hier einführen.}

% wenn von DFA, dann NFA immer complete, hiwneisen (allgemein completeness übertragen)

\todo{auch die konstruktion für kommtuative automateon so handhaben, da
dort kongruenz.}

\todo{diese sichtweise passt in das kongruenzenthema, deswegen auch reachablity
nicht wie rystsov definiert, für die story, irgendwo erwähnen.}

Let $\sim$ be an equivalence relation (induced by any partitioning) on $Q$.
If it is a $\mathcal T_{\mathcal A}$-congruence, then we can construct
the factor automaton and obtain a complete and deterministic automaton.
It turns out that this operation can be performed even if $\sim$
is merely an equivalence relation, but we no longer obtain a deterministic and complete
automaton but an NFA. Formally, $\mathcal A / {\sim}$ \todo{schlau, gleiche notation zu verwenden?}
has as state set the equivalence classes with transition function
$\delta_{\sim} \colon Q / {\sim} \times \Sigma \to \mathcal P(Q / {\sim})$ given by
\todo{notation transition?}
\[
 \applytrans{\delta_{\sim}}{[q]_{\sim}}{x} 
 = \{ [p]_{\sim} \mid p \in \applytrans{\delta}{[q]_{\sim}}{x}\}, 
\]
i.e., if $q \sim q'$, then 
we have a transition labeled by $x$ from $[q]_{\sim}$
to $[\applytrans{\delta}{q}{x}]_{\sim}$
and to $[\applytrans{\delta}{q}{x}]_{\sim}$ in $\mathcal A/{\sim}$.
Note that $\mathcal A /{\sim}$ is a complete NFA (in fact, this construction
can be extended to work for an arbitrary NFA and then preserves completeness,
but as we do not need this generality here, we only apply it to 
deterministic and complete automata).

The next implication follows readily: For every $p, q \in Q$ and $u \in \Sigma^*$ \todo{oder $[\delta] \in \delta()..$
schreiben? aber dann sieht man nicht so leicht, dass nur eine richtung...}
\begin{equation}
\label{constr_sync:eqn:general_factor_transition}
 \delta(p, u) = q \quad\Rightarrow\quad 
 [q]_{\sim} \in \delta_{\sim}([p]_{\sim}, u).
\end{equation}

\todo{referenzen:
%https://en.wikipedia.org/wiki/Edge_contraction
%https://mathworld.wolfram.com/VertexContraction.html
oder das in diskussion?
nicht edge constraction, da keien edge existieren muss.

% in lernagorithmen schonmal gesehen?
%https://www.google.com/search?q=%22state+merging+automaton%22&rlz=1C1GCEU_deDE823DE823&oq=%22state+merging+automaton%22&aqs=chrome..69i57j0i546.6249j1j7&sourceid=chrome&ie=UTF-8

% https://www.google.com/search?q=state+merging+automata+lerning&rlz=1C1GCEU_deDE823DE823&oq=state+merging+automata+lerning&aqs=chrome..69i57j33i10i160l2.4073j1j7&sourceid=chrome&ie=UTF-8
% https://www.researchgate.net/figure/State-merging-approach-for-learning-automata_fig2_288137011
}
We note in passing that the operation corresponds to (repeated) vertex contraction(s)
in the sense of graph theory on the automaton graph while retaining
the labels and not removing  edges, i.e., retaining parallel edges
and self-loops (however, naming it contraction NFA sounds unnatural to us, as in automata theory the term state merging is mostly used in relation to minimization algorithms where
as an equivalence relation the Nerode right-congruence is used).\todo{die nicht definiert, oder nur in diskussion weier oben.}

%
%The construction can also be formalized by interpreting
%by setting 
% closure algorithm auf power automaton und transpose der durch umkehren
% der transitionen entsteht
% interpreation transition von NFA direkt auf potenzautomaten
% 
% und schrittweises verschmezeln von zwei zuständen {p,q}
% defineiren delta({p,q}, x) = menge mit ggfs {p,q}
%
% das in definition\textbf{}

See Figure~\ref{constr_sync:fig:merging_NFA} for an example.

\begin{figure}[ht]
    \centering
    %\begin{minipage}{0.4\textwidth}
    \begin{tikzpicture}
    \node[state] (1) {$1$};
    \node[state, right = of 1] (2) {$2$};
    \node[state, right = of 2] (3) {$3$};
    \path[->] (1) edge [bend left] node {$b$} (2);
    \path[->] (2) edge [bend left] node {$b$} (1)
              (3) edge [loop above] node {$a, b$} (3);
    \path[->] (1) edge [loop above] node {$a$} (1);
    \path[->] (2) edge node {$a$} (3);
    \end{tikzpicture}\hspace{8em}
    %\end{minipage} \hfill
    %\begin{minipage}{0.4\textwidth}
    \begin{tikzpicture}
    \node[state] (1) {$\{1,2\}$};
    \node[state, right = of 1] (2) {$\{3\}$};
    \path[->] (1) edge [loop above] node {$a,b$} (1)
              (1) edge [above] node {$a$} (2)
              (2) edge [loop above] node {$a,b$} (2);
    \end{tikzpicture}
    %\end{minipage} 
    \caption{On the right is the merging NFA for the partitioning of the states
    $\{\{1,2\}, \{3\}\}$ of the automaton on the left.
    }
    \label{constr_sync:fig:merging_NFA}
\end{figure}

\begin{proposition}
\label{constr_sync:prop:D1-directing-implies-sync}
 Let $\mathcal A$ be a deterministic and complete automaton with a trap state $t \in Q$.
 If $\mathcal A_{\text{COND}}$ has a $D_1$-directing word $u$,
 then $u$ is a synchronizing word for $\mathcal A$.
\end{proposition}
\begin{proof}
 We must have $\applytrans{\delta_{\text{COND}}}{[q]_{\simeq}}{u} = \{ [t]_{\simeq} \}$
 for all $q \in Q$.\todo{genauer}
 Let $q \in Q$.
 %Then there exists % das auch zeigen
 %$s \simeq q$ such that $\applytrans{\delta}{s}{u} = t$.
 Consider the state $\applytrans{\delta}{q}{u}$.
 By Equation~\eqref{oben},
 we have $[\applytrans{\delta}{q}{u}]_{\simeq} \in \applytrans{\delta_{\text{COND}}}{[q]_{\simeq}}{u} = \{ \{t\}\}$.
 This implies $\applytrans{\delta}{q}{u} = t$.
\end{proof}

Figure~\ref{constr_sync:fig:merging_NFA} shows that the converse of the statement
from Proposition~\ref{previous} does not hold true, as the automaton on the left
has the synchronizing word $aba$, but its condensation NFA on the right 
has no $D_1$-directing word.
% aber D3-directing hat es, erwähnen?


\todo{solvable synchronizuing iff trap state der von allen erreichbar ist?}

\todo{erwähnen, For deterministic automata a synchronizing word is the same as a $D_1$-directing word}


\subsubsection{The Condensation NFA}


\emph{Restricted} (or \emph{self-loop-determistic}) partially ordered NFA (rpoNFA)
were introduced by Masopust, Kr{\"{o}}tzsch \& Thomazo~\cite{DBLP:journals/iandc/KrotzschMT17}
(see also~\cite{DBLP:journals/lmcs/MasopustK21}). They are defined by the forbidden
pattern depicted in Figure~\ref{constr_sync:fig:rpoNFA}. Formally, a partialy ordered NFA $\semiautomat{A}{Q}{\Sigma}{\delta}$
is a rpoNFA if, for all $q \in Q$ and $x \in \Sigma$, $q \in \applytrans{\delta}{q}{x}$
implies $\applytrans{\delta}{q}{x} = \{q\}$. Expressed in terms of
the reachability relation the condensation NFA is an rpoNFA \todo{acuh allgemein wenn man von NFA started, vielleich es doch so definieren vom start?} 
if and only if $\applytrans{\delta}{q}{x} \not\simeq q$ for $q \in Q$,
then $\applytrans{\delta}{p}{x} \not\simeq q$ for all $p \in Q$ with $p \simeq q$ (or, either of the following two equivalent conditions: (1) $\applytrans{\delta}{[q]_{\simeq}}{x} \cap [q]_{\simeq} \ne \emptyset$
implies $\applytrans{\delta}{[q]_{\simeq}}{x} \subseteq [q]_{\simeq}$,
(2) $\applytrans{\delta}{[q]_{\simeq}}{x} \cap [p]_{\simeq}$ for some $p \not\simeq q$
implies $\applytrans{\delta}{[q]_{\simeq}}{x} \cap [q]_{\simeq} = \emptyset$).


\todo{zustände verkleinern im bild.}
\begin{figure}
    \centering
    \begin{tikzpicture}
    \node[state] (1) {};
    \node[state, right = of 1] (2) {};
    \path[->] (1) edge [loop above] node {$a$} (1)
              (1) edge [above] node {$a$} (2);
    \end{tikzpicture}
    \caption{The forbidden pattern in rpoNFAs.}
    \label{constr_sync:fig:rpoNFA}
\end{figure}


Next, we show that for automata whose condensation NFA is an rpoNFA,
the converse of Proposition~\ref{constr_sync:prop:D1-directing-implies-sync} holds true.






\begin{comment}

\begin{proposition}
 Suppose $\mathcal A$ has a trap state $t\in Q$. % brauche hier nichtmal dass trap state hat?
 Let $S \subseteq Q$ be an upper set and suppose for every $q \in Q$,
 if $\applytrans{\delta}{q}{x} \not\simeq q$,
 then $\applytrans{\delta}{p}{p} \not\simeq q$ \todo{masopust-kratsch eigenschaft für diesen speziellen NFA}
 for every $p \in Q$ with $p \simeq q$.
 Then if $\mathcal A /{\sim_S}$ is sychronizing
 and the condensation NFA of $\mathcal A_{|S}$
 has a $D_1$-directing word, then the condensation NFA of $\mathcal A$
 has a $D_1$-directing word.
\end{proposition}
\begin{proof}
 By assumption there exists $u \in \Sigma^*$
 such that $\applytrans{\delta}{q}{u} \subseteq S$
 for all $q \in Q$.
 %By Equation~\eqref{oben} this implies 
 
 
 % Die Gleichung, die ich hier benutzten möchte gilt nicht! (im Text bemerken?)
 % (sie gilt aber für einzelne zeichen, sonst muss man wörter dazwischenpacken um sich in SCCs zu bewegen)
 %
 % auf beispiel oben
 % [1]aba = {[1], [2]}
 % delta(Q,aba) = {3}, aber 1aba != 3
 %
 % 1aba = 1 geht nicht, existiert also nicht so ein paar wie behauptet!
 Let $[p]_{\simeq} \in \applytrans{\delta_{\text{COND}}}{[q]_{\simeq}}{u}$.
 Then there \todo{genauer}
 exist $\overline q \simeq q$ and 
 $\overline p \simeq p$
 such that $\applytrans{\delta}{\overline q}{u} = \overline p$.
 We have $\overline p \in S$.
 So that $[p]_{\simeq} \subseteq S$ as it is an upper set.
 Furthermore, for each $s \in S$
 there exists $v \in \Sigma^*$
 such that $\applytrans{\delta_{\text{COND}}}{[s]_{\simeq}}{v} = \{ [t]_{\simeq} \}$.
 The latter implies $\applytrans{\delta_{\text{COND}}}{[s]_{\simeq}}{uv} = \{ [t]_{\simeq} \}$
 for each $s \in S$.
\end{proof}
\end{comment} 

\todo{topological ordering/sorting definieren}

\todo{nach vorne verlegen.}
\begin{lemma}
\label{prelim:lem:sync_und_trap_state}
 If $\semiautomat{A}{Q}{\Sigma}{\delta}$
 has a trap state $t \in Q$ and is synchronizing,
 then the set of synchronizing states is $\{t\}$
 (the maximal strongly connected component in the induced partial order).\todo{sprechweise induced partial order einführen.}
\end{lemma}

\todo{hier muss man genau unterscheiden welchen automateontyp man betrachtet}

\begin{lemma}
\label{constr_sync:lem:rpoNFA_path-faithfulness}
 Let $\semiautomat{A}{Q}{\Sigma}{\delta}$ be a DFA.
 Suppose $\applytrans{\delta}{[q]_{\simeq}}{u} = \{ p \}$. % weil keine zustände in wzischen-SCC's eingeführt,
 % wo es loops kann C_1 -> C_2 -> {p}, C_2 noch loopzustände, beispiel dazu
 Then $\applytrans{\delta_{\text{COND}}}{[q]_{\simeq}}{u} = \{ [p]_{\simeq} \}$.
\end{lemma}
\begin{proof}\todo{index simeq dranschreiben}
 For $x \in \Sigma$ the property is clear.\todo{definition transition genauer, gilt für beliebig condNFAs}
 Consider $u = xv$ and set $T = \applytrans{\delta_{\text{COND}}}{[q]}{x} = $ \todo{wie oben def}
 % v = eps, also für eps als IA
 Inductively, for each $[r] \in T$
 %
 % r nehmen, muss aber nicht auf p landen! also kann man IV nicht anwenden,
 % aber für sync-wort gilt es ja!
 
 
\end{proof}


\todo{anderes argument: wenn $u$ D3-directing in rpoNFA mit sink state,
dann $u^i$ für ein $i$ $D1$-directing. male dreieck, wenn $q$
in $t$ und $p$ endet, von $p$ auf $t$, betrachte $uu$.}



\todo{mal DFA, mal complete and det. automat -> vereinheitlichen, PDFA, NFA}
\begin{proposition}
 Let $\semiautomat{A}{Q}{\Sigma}{\delta}$ be a DFA having a trap state $t \in Q$
 and whose condensation NFA is a rpoNFA.
 If $u \in \Sigma^*$ is a synchronizing word for $\mathcal A$,
 then $u$ is a $D_1$-directing word in the condensation NFA.
\end{proposition}
\begin{proof}
 Let $[q_1], \ldots, [q_m]$ be a topological ordering.
 Define upper sets $S_i = [q_i] \cup \ldots \cup [q_m]$.
 We show inductively (over $i$) that if $\mathcal A_{|S_i}$
 has a synchronizing word $v$, then $v$ is $D_1$-directing in its condensation NFA.
 
 
 By Lemma~\ref{prelim:lem:sync_und_trap_state}, 
 we have $S_m = \{t\}$.
 If $i = m$, then $\mathcal A_{|S_m}$ is an automaton with a single state,
 it equals its condensation NFA and $\Sigma^*$ is the set of synchronizing
 words as well as the set of $D_1$-directing words.
 So the statement is true.
 
 
 
 Suppose $1 \le i < m$ and $u$ synchronizes $\mathcal A_{|S_i}$.
 Then we can write $u = vxw$ with $x \in \Sigma$ and $v,w \in \Sigma^*$
 such that $\applytrans{\delta_{|S_i}}{[q_i]_{\simeq}}{vx} \cap S_{i+1} \ne \emptyset$
 and $\applytrans{\delta_{|S_i}}{[q_i]_{\simeq}}{vx} \cap S_{i+1} = \emptyset$.
 As the condensation NFA is a rpoNFA, 
 if $\applytrans{\delta_{|S_i}}{[q_i]_{\simeq}}{vx} \cap S_{i+1} \ne \emptyset$,
 then $\applytrans{\delta_{|S_i}}{[q_i]_{\simeq}}{vx} \subseteq S_{i+1}$.
 The word $vxw$ synchronizes $\mathcal A_{|S_{i+1}}$.
 Inductively, we find that it is a $D_1$-directing word 
 for the condensation NFA of $\mathcal A_{|S_{i+1}}$.
 This yields $\applytrans{\delta_{|S_{i+1}}}{}{}$
 
 
 
 So, $\applytrans{\delta_{|S_i}}{}{}$
 
 
 
 Then it also synchronizes $\mathcal A_{|S_{i+1}}$.
 By induction hypothesis, the condensation NFA of $\mathcal A_{|S_{i+1}}$
 has a $D_1$-directing word $u \in \Sigma^*$.
 There  \todo{genauer}
 exists $x \in \Sigma$ such that $\applytrans{\delta}{q}{x} \in S_{i+1}$ %masopust-kratsch eigenschaft & sync
 for each $q \in Q$ with $q \simeq q_i$.
 By the defining equation of the transition \todo{auf gleichung verweise.}
 we have $\applytrans{\delta_{\text{COND}}}{[q_i]_{\simeq}}{x} \subseteq \{ [s]_{\simeq} \mid s \in S_{i+1} \}$.
 As $\applytrans{\delta_{\text{COND}}}{[s]_{\simeq}}{u} = \{ [t]_{\simeq} \}$
 and $\applytrans{\delta_{\text{COND}}}{[s]_{\simeq}}{xu} = \{ [t]_{\simeq} \}$, \todo{genauer}
 it follows that $xu$ is $D_1$ directing for the condensation NFA of $\mathcal A_{|S_i}$
 (which is the condensation NFA of $\mathcal A_{|S_{i+1}}$
 with the additional state $[q_i]_{\simeq}$ and appropriate transitions).
\end{proof}






















\todo{semi-DFA nicht extra sagen. das auch noch an anfang irgendwo}


\begin{comment} 
A DFA $\semiautomat{A}{Q}{\Sigma}{\delta}$ % t trap state, s sync state
is \emph{residual synchronizing} if it is synchronizing
and for every synchronizing word $u \in \Sigma^*$ with $\applytrans{\delta}{Q}{u} = \{s\}$,
if $u = vw$ with $v,w \in \Sigma^*$
and $\applytrans{\delta}{q}{v} \simeq p$ for $p, q \in Q$,
then $\applytrans{\delta}{p}{w} = s$. 
\todo{siehe kommentar im quelltext
% ein wort ist residual-sync, wenn es dies erfüllt, 
% gibt stets residual-sync wörter durch lange genugt machen (u^{|Q|}), definition ist, dass es alle sind.
% damit constraint residual synchronization problem auf NFA reduzierbar
%
% kürzestes residual-sync wort?
%
% wenn Syn(A) unter u^{-1}L abgeschlossen, dann residual-sync automat.
%
% sinnvolle beispiele für residual-sync automaten? für SCC heisst es, jedes suffix ist sync, und
% das gilt natürlich nicht, also eher für pto automaten interessant
%
% kommutative automaten sind residual-sync?
}
\end{comment}

A DFA $\semiautomat{A}{Q}{\Sigma}{\delta}$ % t trap state, s sync state
is \emph{residual synchronizing} if it is synchronizing
and for every synchronizing word $u \in \Sigma^*$ with $\applytrans{\delta}{Q}{u} = \{s\}$,
if $u = vw$ with $v,w \in \Sigma^*$
and $q \not\simeq \applytrans{\delta}{q}{v}$ for $p, q \in Q$,
then $\applytrans{\delta}{p}{w} = s$. % also wenn praktsich in größerer SCC-komponente landet
%
% ne, macht keinen sinn.
% will dass alle SCCs hinter dem von delta(q,v) in top sort von w gesynct werden.
%
% oder brauche ich nicht sowas wie, dass w konfluent ist für die, die auseinandergehen?















\subsubsection{Synchronizing Words with Respect to a Constraint Language}

If the partitioning is given by the equivalence relation $\simeq$
induced by the reachability order (see Section~\ref{reachability}),
then we call the merging NFA the \emph{condensation NFA}, denoted by 
$\semiautomat{A_{\text{COND}}}{Q/{\simeq}}{\Sigma}{\delta_{\text{COND}}}$.
It corresponds to collapsing, or merging, the strongly connected\todo{refernzen?}
components of the automaton (this operation is called condensation
in graph theory, and the condensation NFA corresponds to the condensation
of the automaton graph).

\begin{proposition} \todo{nicht unbedingt solvable! diskutieren, auch mit pspace konstruktionen vergleichen die ja auch trap states haben}
 Let $\mathcal A$ be a complete and deterministic automaton with a trap state $t \in Q$ and $u \in \Sigma^*$.
 The word $u$ synchronizes $\mathcal A$
 if and only if it is $D_1$-directing in $\mathcal A_{\text{COND}}$.
\end{proposition}
\begin{proof}
 The trap state must be the synchronizing state in $\mathcal A$.
 Let $u$ be a synchronizing word for $\mathcal A$.
 Then $\applytrans{\delta}{Q}{u} = \{t\}$ and 
 the defining equation of the transition function \todo{mit nummer zitieren}
 gives, for each $q \in Q$,
 \[ 
  \applytrans{\delta_{\text{COND}}}{[q]_{\simeq}}{u}
  = \{ [p]_{\simeq} \mid p \in \applytrans{\delta}{[q]_{\simeq}}{u} \} 
  = \{ [p]_{\simeq} \mid p \in \{ t \} \}
  = \{ [t]_{\simeq} \}.
 \]
 So the word is $D_1$-directing for $\mathcal A_{\text{COND}}$.
 \todo{das ist falsch, da man in jedem schritt eine art closure operation ausführen muss!
 obiges beispiel zeigt es!}
 
 Conversely, suppose $u$ is a $D_1$-directing word for $\mathcal A_{\text{COND}}$.
 
 % ne, im Beispiel ab Gegenbeispiel!!!!
 
\end{proof}




\todo{ganz zu anfang let $A = $ schreiben, in dieser section}

\todo{mit rystsov definitionen vergleichen. auch recahblity order und folgerungen damit; aber hier kongruenzperspektive}
The automaton $\mathcal A$ is called \emph{solvable}, if
there exists a topological sorting 
\[ 
[q_1]_{\simeq}, [q_2]_{\simeq},
\ldots, [q_m]_{\simeq}
\] such that $\mathcal A_{|S_m}$
is commutative and for all $i \in \{1,\ldots,m-1\}$
the Rees factor automata $( \mathcal A_{|S_i} ) / {\sim_{S_{i+1}}}$
are commutative,
where $S_i = [q_i]_{\simeq} \cup [q_{i+1}]_{\simeq} \cup \ldots \cup [q_m]_{\simeq}$ 
(note that these are upper sets).
Actually, if the condition holds true for some topological ordering,
then it holds true for all topological sortings, as
the factor automata arise out of the strongly connected components
and an additional trap state.\todo{genauer? jordan-hölder in rystsov erwähnen?}

\todo{den struktursatz auch nochmal hinschreiben? permutation reset automata, und im Pol(GComm)?
verweis auf Schur lemma}

A synchronizing word for $\mathcal A$ synchronizes $\mathcal A_{|S_m}$ as well.
As it is commutative, by Fact~\ref{comm_sync_state_is_trap_state}
and as it is strongly connected, we conclude $S_m = \{q_m\}$
and $q_m$ is a trap state.

\begin{lemma}
\label{constr_sync:lem:solvable_sync_trap_state}
 If $\mathcal A$ is solvable and synchronizing, then the synchronizing
 state is a unique trap state.
\end{lemma}

Let $q \in S_i$ and $x \in \Sigma$. If $\applytrans{\delta}{q}{x} \sim q$,
then $\applytrans{\delta}{p}{x} \sim q$ for all $p \in S_i$. \todo{so oben auch formulieren?}
This follows \todo{die facts dazu ziteiren?}
as in a commutative automaton every letter either permutes a strongly connected
components or maps every state of a strongly connected component into states
outside of this strongly connected component.

\begin{lemma}
 Let $\mathcal A$ be solvable, $q \in Q$ and $x \in \Sigma$.
 If $\applytrans{\delta}{q}{x} \sim q$,
 then $\applytrans{\delta}{p}{x} \sim q$
 for all $p \in Q$ with $p \sim q$.
 %$p, q \in Q$ and $p \sim q$,
 %then for each $x \in \Sigma$
 %we either have $\applytrans{\delta}{p}{x} \sim \applytrans{\delta}{q}{x}$
 %for 
\end{lemma}

\todo{masopust, kratsch bild, pattern}

\todo{$D_i$-directing einführen und diskutieren, dass nur D1 geht.}























\todo{entscheidungsverfahren für constr-sync nur diskutieren und
auf paper verweisen.}

\todo{- pto => monotonic aber nicht umgekehrt (aus sync-games paper beispiel)}

\todo{auskommentierten teil später noch irgendwie einbauen.}
\begin{comment} 

Let $\mathcal A$ be an automaton. An upper set $S \subseteq Q$
is \emph{maximal} if $S \ne Q$ and for every upper set $S \subseteq T \subseteq Q$
we have $S = T$ or $T = Q$. %composition series defineren


% composition series erwähnen und kommentieren entspricht top sort der SCC's

% S maximal in T

A \emph{composition series} is a sequence of upper
set 
\[
 S_0 \subseteq 
\]


An automaton is solvable %mit composition series



An automaton is solvable if for each $q \in Q$
with $\Gamma = \{ x \in \Sigma \mid \applytrans{\delta}{q}{x} \cap [q]_{\sim} \ne \emptyset \}$
the automaton $\mathcal A_{[q]_{\sim}, \Gamma}$
is commutative. % äquiavlent permutation autoamta

\todo{ 
aber [q]_{\sim} kein upper set, aber defintiion geht für die gamma abgeschlossen; dafür müsste diese definition oben angepasst werden.

und das geht dann auch nur für solvable, also eher schlecht damit zu definieren,
oder das mit in defiition packen.

}


\begin{proposition}
 The automaton $\mathcal A$ is solvable if and only if
 %every composition factor is a permutation reset automaton.
 for every upper set $S \subseteq Q$
 and maximal upper set $T$ in $\mathcal A / {\sim_S}$
 the Rees quotient for $T$ in $\mathcal A / {\sim_S}$ \todo{oder Rees factor?}
 is a commutative permutation-reset automaton.
\end{proposition}







% wenn sync maximales SCC [q] = {q} irgendwo hinschreiben?




% die amsopust-kratsch form wichtig für D_3
% und übrigens keine undefinierten übergänge (also complete)

\todo{auf relation in $\Sigma^*$ erweitern.}

Clearly, as we are collapsing states, % mit funktion beschreiben
\begin{equation}\label{constr_sync:eqn:cond_paths} 
\applytrans{\delta}{s}{u} = t \Rightarrow ([s]_{\sim}, u, [t]_{\sim}) \in \delta_{\text{COND}}
\end{equation}
(this holds in general for collapsing, or merging,
of states).\todo{genauer, induktion, mit ziwe zuständen starten.}
 If there exists $([s]_{\sim}, u, [t]_{\sim}) \in \delta_{\text{COND}}$,
 %
 % ne, wenn man über collapsed geht nicht für alle
 then there exists $\overline s \sim s$
 and $\overline t \sim t$
 such that $\applytrans{\delta}{\overline s}{u} = \overline t$. \todo{beweisen.}
\end{comment}
 



\todo{oder war $D_1$ alles in einen?}


\begin{proposition} \todo{solvable DFA? / FA / auotmaton}
 Let $\semiautomat{A}{Q}{\Sigma}{\delta}$ be solvable with synchronizing state $t \in Q$
 and $u \in \Sigma^*$.
 The word $u$ synchronizes~$\mathcal A$, i.e. $\applytrans{\delta}{Q}{u} = \{t\}$,
 if and only if it is $D_1$-directing for the condensation NEA~$\mathcal A_{\text{COND}}$,
 in which case $\applytrans{\delta_{\text{COND}}}{Q/{\simeq}}{u} = \{ [t]_{\simeq} \}$.
\end{proposition}
\begin{proof}
 First, suppose $u$ synchronizes $\mathcal A$ with synchronizing state $q$.
 By Equation~\eqref{constr_sync:eqn:cond_paths},
 for every $s \in Q$
 we have $([s]_{\sim}, u, [q]_{\sim}) \in \delta_{\text{COND}}$.
 Now assume $([t]_{\sim}, u, [p]_{\sim}) \in \delta_{\text{COND}}$,
 then with the above $[p]_{\sim} = [q]_{\sim}$. \todo{genauer}
 
 \todo{sprechweise to the state}
 Now, suppose $u$ is $D_1$-directing for $\mathcal A_{\text{COND}}$
 to the state $[q]_{\sim}$. % ist letzter in A, also = {q}
 
 
 Let $s \in Q$.
 Then $([s]_{\sim}, u, [q]_{\sim}) \in \delta_{\text{COND}}$.
 Then % obige erklärung
 there exists $\overline s \sim s$
 such that $\applytrans{\delta}{\overline s}{u} = q$. % hier = {q} genutzt! als lemma?
 
 
 
 
 %
 % induktion über topsort, und da entweder self-loop oder beide weiter
 % geht es
 
 
 Suppose $\applytrans{\delta}{s}{u} = p$,
 then with Equation~\eqref{todo} \todo{genauer}
 $p = q$. % da sync
 
 %
 % aber dann habe ich solvable gar nicht gebraucht!!!!!
 %
 % häöchstens in der eigenschaft dass defineirt, also complete!
 %
 % dann wären complete NEA ja immer in NP! also missing transition
 % für pspace!
 
 \todo{$\simeq$ nicht $\sim$!}
 
 \todo{$q_m = t$}
 \todo{vielleicht diese aussage als separates lemma?}
 \todo{erwähnen daraus composition series durch vereinigungen qi bis qm?}
 Let $[q_1]_{\sim}, [q_2]_{\sim}, \ldots, [q_m]_{\sim}$
 be a topological sorting in the induced partial order of the reachability order. \todo{zeigen, dass es letzer zustand sein muss. und solvable dann trap state (gleich menge sync states singleton)}
 We show by induction over $i \in \{1,\ldots,m\}$ the following:
 \todo{einfacher formulieren $q = q_i$ setzen}
 If $\applytrans{\delta_{\text{COND}}}{[q_i]_{\sim}}{w} = \{ [q_m]_{\sim} \}$ 
 for some $i \in \{1,\ldots,m\}$ and $w \in \Sigma^*$,
 then $\applytrans{\delta}{q}{w} = q_m$
 for all $q \in Q$ with $q \sim q_i$.
 If $q \sim q_m$, then $q = q_m$
 and with Equation~\eqref{constr_sync:eqn:cond_paths} \todo{bzw, weil der trap state}
 we find $\applytrans{\delta}{q}{w} = q_m$.
 So, suppose $1 \le i < m$.
 Then $w \ne \varepsilon$ and 
 we can write $w = uxv$ with $u,v \in \Sigma^*$ and $x \in \Sigma$
 such that $\applytrans{\delta}{q_i}{u} \sim q_i$
 and $\applytrans{\delta}{q_i}{ux} \not\sim q_i$.
 As $\mathcal A$ is solvable, if $q \sim q_i$ \todo{lemma zitieren?}
 then $\applytrans{\delta}{q}{ux} \not\sim q_i$ as well.
 Consider the state $s = \applytrans{\delta}{q}{uxv}$.
 By Equation~\eqref{todo}, we have
 \[
  % [s]_{\simeq} \in \applytrans{\delta_{\text{COND}}}{[\applytrans{\delta}{q}{ux}]_{\simeq}}{v}.
  [s]_{\simeq} \in \applytrans{\delta_{\text{COND}}}{[q]_{\simeq}}{w}. %[q] = [q_i]
 \]
 So $[s]_{\simeq} = [q_m]_{\simeq}$, as we only end up in $\{ [q_m]_{\simeq}\}$,
 which gives $s = t$. % wieder nur completeness....
 
 
 
 
 
 
 
 Let $i < j \le m$ with $q_j \sim \applytrans{\delta}{q}{ux}$.
 We have\todo{genauer?}
 \[ 
  \applytrans{\delta}{}{}
 \] 
 
 
 
 
 
 
 
 
 
 
 
 
 
  Let $q \in Q$
 and suppose $q \sim q_i$ for some $i \in \{1,\ldots,m\}$. 
 %We have $\applytrans{\delta}{[q]_{\sim}}{u} = \{ [q_m]_{\sim} \}$.
 If $i = m$, then 
 as $[q_m]_{\sim} = \{q_m\}$ we have $q = q_m$
 and $\applytrans{\delta}{q_m}{u} = q_m$ follows.
 Otherwise, suppose $1 \le i < m$.
 \todo{die gleichung einmal am anfang hinschreiben.}
 As $\applytrans{\delta_{\text{COND}}}{[q]_{\sim}}{u} = [q_m]_{\sim}$,
 there exist $p \in Q$
 with $q \sim p$ and $\applytrans{\delta}{p}{u} = q_m$.
 We have $u \ne \varepsilon$ and can write $u = vxw$ with $v,w \in \Sigma^*$
 and $x \in \Sigma$
 such that $\applytrans{\delta}{p}{v} \sim p$
 and $\applytrans{\delta}{p}{vx} \not\sim p$.
 As $\mathcal A$ is solvable, \todo{genauer?}
 this implies $\applytrans{\delta}{q}{vx} \not\sim p$.
 We have \todo{diese gleichung oben erwähnen. und auch warum nur teilmengenbeziehung}
 \[
  [q_m]_{\sim} = \applytrans{\delta_{\text{COND}}}{[q]_{\sim}}{u}
  \supseteq \applytrans{\delta_{\text{COND}}}{[\applytrans{\delta}{q}{vx}]_{\sim}}{w}
 \]
 and the latter set is non-empty as the automaton is complete.
 So, as $[q_m]_{\sim} = \{q_m\}$ we have in fact equality in the above equation.
 
 
 
 
 
 
 
 
 
 
 
 
 
 
 
 
 
 
 
\end{proof}

\todo{subset probleme und schranken nicht in diss, dafür originalarbieten,
aber sonst besser in diss}



% top sort, dann induktion q_i nach {q_{i+1}, ..., q_n

\todo{in regulärer sprache!}



 \todo{bei nea einführen sagen auf funktion nach postenzmenge, und kurz automateongraphen ansprechen.}
 
 
\begin{proposition}
 Let $\automat{B}{P}{\Sigma}{\mu}{p_0}{F}$
 be a PDFA.
 A partially ordered
 $n$-state
 $D_1$-directable NFA $\semiautomat{A}{Q}{\Sigma}{\delta}$ has a $D_1$-directing word 
 in $L(\mathcal B)$ of length at most $(n - 1) \cdot |P|$. 
 \todo{im text erwähnen easily seen to be sharp even for deterministic and complete,
 folgerung damit auch synchronizing word vopn partially orderer DFA}
\end{proposition}
\begin{proof}
 Let $Q = \{q_1, q_2, \ldots, q_n\}$
 be a topological sorting of the states.\todo{genauer hinschreiben mit reachability was da heißt}
\begin{comment} 
 As we have a single state to which every other state is mapped\todo{englisch, is mapped to?}
 to, this single state must be the last state~$q_n$ in the topological 
 ordering. \todo{genauer?}
 Now, let $i \in \{1,\ldots,n-1\}$
 and consider $q_i \in Q$.
 Then there must exist $x_i \in \Sigma$ and  $s \in \{q_{i+1}, \ldots, q_n\}$
 such that $(q_i, x_i, s) \in \delta$, otherwise
 this state cannot be mapped to $q_n$. \todo{da es nur vorwärts geht, genauer?}
 For $x \in \Sigma$ and $t \in Q$
 with $(q_i, s, t) \in \delta$
 we have $t \in \{q_i, q_{i+1}, \ldots, q_n\}$
 as $\mathcal A$ is partially ordered,
 i.e., the letter either induces a self-loop or moves forward in the partial order.
 So, for the word $u = x_1 x_2 \cdots x_n$
 we \todo{nehme ich jetzt die relationenschreibweise oder funktion in potenzmenge?}
 have $q_n \in \applytrans{\delta}{Q}{u}$.
 
 
 % praktisch blockcode Sigma^{|P|}
\end{comment}
 Let $u \in L(\mathcal B)$ be a shortest $D_1$-directing word. 
 Write $u = u_1 u_2 \cdots u_{n-1}$
 where $u_i \in \Sigma^*$ for $i \in \{1,\ldots,n-1\}$
 is shortest possible
 such that $\applytrans{\delta}{q_i}{u_1 \cdots u_i} \subseteq \{q_{i+1}, \ldots, q_n\}$
 (note that $u_i = \varepsilon$ is possible).
 By choice, if $u_i\ne \varepsilon$,
 we can write $u_i = v_i x_i$ with $x_i \in \Sigma$ and $v_i \in \Sigma^*$
 such that $\applytrans{\delta}{q_i}{v_i} = \{q_i\}$ \todo{bemerken muss singleton sein wegen pto}
 and $\applytrans{\delta}{q_i}{x_i} \subseteq \{q_{i+1}, \ldots, q_n\}$.
 In this case, $\applytrans{\delta}{q_i}{v} = \{q_i\}$ \todo{prefix/suffix/factor/infix definiert?}
 for each prefix $v$ of $u_1 \cdots u_{i-1}$.
 Note that $\applytrans{\delta}{\{q_1, \ldots, q_i\}}{u_1 \cdots u_i}
 \subseteq \{q_{i+1}, \ldots, q_n\}$ and as $u$
 is $D_1$-directing this implies \todo{problematischer schritt, da oben nur teilmenge, also muss es
 ja nicht alles synchronisieren. ist wie diese residual-eigenschaft}
 \begin{equation}
 \label{constr_sync:eqn:pto_NFA}
     \applytrans{\delta}{\{q_{i+1}, \ldots, q_n\}}{u_{i+1} \cdots u_{n-1}}
      = \{q_n\}.
 \end{equation}
 Assume $|v_i| \ge |P|$. Then $v_i$ induces a loop in $\mathcal B$
 and we find a proper prefix $v_i' \in \Sigma^*$ with $|v_i'| < |v_i|$ 
 and \todo{genauer?}
 $\applytrans{\mu}{p_0}{u_1 \cdots u_{i-1} v_i} = \applytrans{\mu}{p_0}{u_1 \cdots u_{i-1} v_i'}$.
 %As $\applytrans{\delta}{q_i}{v_i'} = \{q_i\}$
 %and $\applytrans{\delta}{q_i}{x_i} \subseteq \{q_{i+1},\ldots,q_n\}$
 We have $\applytrans{\delta}{q_i}{u_1 \cdots u_{i-1} v_i'x_i}
 \subseteq \{ q_{i+1}, \ldots, q_n \}$.
 Furthermore,
 $u_1 \cdots u_{i-1} v_i' x_i u_{i+1} \cdots u_n \in L(\mathcal B)$
 and with Equation~\eqref{constr_sync:eqn:pto_NFA},
 \[
 \applytrans{\delta}{\{q_1, \ldots, q_n\}}{u_1 \cdots u_{i-1} v_i' x_i u_{i+1} \cdots u_n}
  = \applytrans{\delta}{\{q_{i+1},\ldots,q_n\}}{ u_{i+1} \cdots u_n }
  = \{q_n\}
 \]
 so that this word is a shorter $D_3$-directing
 word, which is excluded by assumption.
 Hence $|u_i| \le |P|$ for all $ \in \{1,\ldots,n\}$.
 This yields $|u| \le (n-1)\cdot |P|$.
 % problem verändert teil dahinter/davor
 % neues wort auch in L(B)
%  ...
%  such that %$\applytrans{\delta}{q_i}{u_i} \subseteq \{q_{i+1}, \ldots, q_n\}$
%  $\applytrans{\delta}{q_i}{u_i u_{i+1} \cdots u_m} = \{q_n\}$.
%  Note that this implies $\applytrans{\delta}{\{q_i, q_{i+1}, \ldots, q_n\}}{u_i u_{i+1} \cdots u_n} = \{q_n\}$ and that $u_i = \varepsilon$ is possible. \todo{genauer}
 %
 % letztens zeichen bewegt, und man kann kürzen und inducitkv annehmen alle danach <|P|
 %
 % wo ist oben das problem?
 % 
 % we show |u_i| < |P| induktively
 %
 % komplette wort u_1 .... v' .... u_m noch sync
 % ne, ich nehme mir ja nur blöcke aus dem wort
\end{proof}



%%%%%%%%%%%%%%%%%%%%%%%%%%%%%%%%%%%%%%%%%%%%%%%%%%%%%%%%%%%%%%%%%%%%%%%%%%%%%%%%%%%%%%
% NEU



\todo[inline]{

Gleiche Problem wie vorher, dass kürzere $v_i'$ kann in $[q_i]$
drin rumpermutieren. $x_i$ mappt zwar raus, aber das ergebnis muss von $u_{i+1} \cdots u_{m-1}$
nicht mehr synchronisiert werden. (mit zyklenargument?)

}
\begin{proposition}
 Let $\automat{B}{P}{\Sigma}{\mu}{p_0}{F}$
 be a PDFA and $\semiautomat{A}{Q}{\Sigma}{\delta}$
 a solvable automaton.
 If $\mathcal A$ is synchronizing,
 it has a synchronizing word of length at most $(n - 1)\cdot |P|$.
\end{proposition}
\begin{proof}
 Let $[q_1]_{\simeq}, [q_2]_{\simeq}, \ldots, [q_m]_{\simeq}$
 be a topological sorting of the strongly connected components.\todo{erwähnen induced order?}
 Suppose $u \in L(\mathcal B)$ is a shortest synchronizing word for $\mathcal A$.
 Write $u = u_1 u_2 \cdots u_{m-1}$ where $u_i \in \Sigma^*$ \todo{verträglichkeit der ordnung bzw, dass es nie zurückgeht erwähnen?}
 for $i \in \{1,\ldots,m-1\}$ is shortest possible such that \todo{indiczes simeq setzen}
 $\applytrans{\delta}{[q_i]}{u_1 \cdots u_i} 
 \subseteq [q_{i+1}]_{\simeq} \cup \ldots \cup [q_m]_{\simeq}$
 (note that $u_i = \varepsilon$ is possible).
 If $u_i \ne \varepsilon$, then we
 can write $u_i = v_i x_i$ for some $x_i \in \Sigma$
 and $v_i \in \Sigma^*$.
 Note that this implies 
 $\applytrans{\delta}{[q_1] \cup \ldots \cup [q_i]}{u_1 \cdots u_i}
 \subseteq [q_{i+1}]_{\simeq} \cup \ldots \cup [q_m]_{\simeq}$.
 By choice, we have $\applytrans{\delta}{[q_i]}{v_i} \cap [q_i] \ne \emptyset$.
%  However, with Lemma~\ref{immer-komplett-raus},
%  this implies   \todo{genauer? sogar gleich da endlich, permutiert}
% %  \[ \applytrans{\delta}{[q_i]}{v} \subseteq [q_i]
% %  \]
% %  for every prefix $v$ of $v_i$.
 So $\applytrans{\delta}{[q_i]}{x_i} \not\subseteq [q_i]$, and with Lemma~\ref{todo}, 
 \begin{equation}
     \applytrans{\delta}{[q_i]}{x_i} \subseteq [q_{i+1}] \cup \ldots [q_m],
 \end{equation}
 i.e., all transition labeled with $x_i$ and starting at a state in $[q_i]$
 point to states outside of $[q_i]$. %kann man auch schlussfolgern aus der Tatsache , das v_i [q_i] permutiert.
 Assume $|v_i| \ge |P|$. Then $v_i$ induces a loop in $\mathcal B$
 and we find a proper prefix $v_i' \in \Sigma^*$ with $|v_i'| < |v_i|$ 
 and \todo{genauer?}
 $\applytrans{\mu}{p_0}{u_1 \cdots u_{i-1} v_i} = \applytrans{\mu}{p_0}{u_1 \cdots u_{i-1} v_i'}$.
 %As $\applytrans{\delta}{q_i}{v_i'} = \{q_i\}$
 %and $\applytrans{\delta}{q_i}{x_i} \subseteq \{q_{i+1},\ldots,q_n\}$
 So, $u_1 \cdots u_{i-1} v_i' x_i u_{i+1} \cdots u_n \in L(\mathcal B)$.
 With Equation~\eqref{qivi_in_qi},
 \begin{align*}
   \applytrans{\delta}{Q}{u_1 \cdots u_{i-1} v_i' x_i u_{i+1} \cdots u_{m-1}} 
    & \subseteq \applytrans{\delta}{[q_i] \cup \ldots \cup [q_m]}{v_i'x_i u_{i+1} \cdots u_{m-1}} \\
    & \subseteq \applytrans{\delta}{[q_i] \cup \ldots \cup [q_m]}{x_i u_{i+1}\cdots u_{m-1}} \\
    & \subseteq \applytrans{\delta}{[q_{i+1}] \cup \ldots \cup [q_m]}{u_{i+1}\ldots u_{m-1}}
 \end{align*}
 
 
 We have $\applytrans{\delta}{[q_i]}{u_1 \cdots u_{i-1} v_i'x_i}
 \subseteq [q_{i+1}] \cup \ldots \cup [q_n]$.
 Furthermore,
 
 and \todo{auf obige gleichung verweisen?}
 \begin{multline*}
     \applytrans{\delta}{Q}{u_1 \cdots u_{i-1} v_i' x_i u_{i+1} \cdots u_{m-1}}
    =  \applytrans{\delta}{[q_i] \cup \ldots \cup [q_m]}{v_i' x_i u_{i+1} \cdots u_{m-1}}
 \end{multline*}
 
 
 
 \[
 \applytrans{\delta}{\{q_1, \ldots, q_n\}}{u_1 \cdots u_{i-1} v_i' x_i u_{i+1} \cdots u_n}
  = \applytrans{\delta}{\{q_{i+1},\ldots,q_n\}}{ u_{i+1} \cdots u_n }
  = \{q_n\}
 \]
 so that this word is a shorter $D_3$-directing
 word, which is excluded by assumption.
 Hence $|u_i| \le |P|$ for all $ \in \{1,\ldots,n\}$.
 This yields $|u| \le (n-1)\cdot |P|$.
 
 
 
\end{proof}